% 【本文件写什么】impl 实现层：virtio-mmio
% - 定位：VirtIO MMIO 网卡，RISC-V。
% - crate 路径：os/components/wateros-driver/driver-network/network-impl/impl-virtio-mmio/src/
% - 如何实现对应 api-v0 trait：关键类型、状态机、数据结构
% - 算法或硬件交互流程，可用伪代码/流程图
% - 本 impl 启用的 Cargo [features] 与 arch feature，impl-riscv64 等
% - 与其它 impl 变体的差异、切换 feature 时的影响
% - 性能/内存假设与已知限制
% 【事实来源】
% - 上述 crate 源码与 Cargo.toml
% - os/feature-tree.txt 中指向本 impl 的 feature 链

\subsubsection{impl-virtio-mmio，network}

\paragraph{启用与探测。}由\code{impl-qemu-riscv64-opensbi}打开。平台在\code{virtio,mmio}节点且\code{device\_id==1}时调用\code{VirtioNetDevice::from\_mmio}，成功后\code{register\_network\_device}并记录MAC。

\paragraph{实现要点。}\code{VirtIONet}包装于\code{VirtioNetDevice}；\code{VirtioMmioHal}与块设备共用帧分配器DMA策略，恒等映射、连续PPN校验。接收缓冲\code{RX\_BUF\_LEN==2048}，满足virtio net header与以太网MTU下限。\code{send}将完整以太网帧提交传输队列；\code{receive}从设备拉取帧写入调用方缓冲，无数据返回0。

\begin{rustcode}
pub struct VirtioNetDevice {
    inner: VirtIONet<VirtioMmioHal, MmioTransport<'static>>,
}

impl NetworkDevice for VirtioNetDevice {
    fn mac_address(&self) -> [u8; 6] {
        self.inner.mac_address()
    }
    fn send(&mut self, buf: &[u8]) -> DriverResult<()> { /* ... */ }
    fn receive(&mut self, buf: &mut [u8]) -> DriverResult<usize> { /* ... */ }
}
\end{rustcode}

\paragraph{限制。}与块MMIO路径相同：未使用\code{DeviceInfo.irq}，收包依赖smoltcp轮询\code{receive}。多网卡时仅首张被\code{stack::init}选用，其余仅存在于全局表与devfs诊断信息中。
